\documentclass[14 pt]{extarticle}

	\usepackage[frenchb]{babel}
	\usepackage[utf8]{inputenc}  
	\usepackage[T1]{fontenc}
	\usepackage{amssymb}
	\usepackage[mathscr]{euscript}
	\usepackage{stmaryrd}
	\usepackage{amsmath}
	\usepackage{tikz}
	\usepackage[all,cmtip]{xy}
	\usepackage{amsthm}
	\usepackage{varioref}
	\usepackage{geometry}
	\geometry{a4paper}
	\usepackage{lmodern}
	\usepackage{hyperref}
	\usepackage{array}
	 \usepackage{fancyhdr}
	 \usepackage{float}
\renewcommand{\theenumi}{\alph{enumi})}
	\pagestyle{fancy}
	\theoremstyle{plain}
	\fancyfoot[C]{} 
	\fancyhead[L]{Interrogation}
	\fancyhead[R]{2025-2026}\geometry{
 a4paper,
 total={170mm,257mm},
 left=20mm,
 top=20mm,
 }
	
	
	\title{Interrogation chapitre 3}
	\date{}
	\begin{document}

\begin{center}{\Large Interrogation chapitre 3}\\ 
 \end{center}
 
 
 \subsection*{Exercice 1 (5 points)}
 
 Effectuer les calculs suivants en détaillant les étapes : 
 \begin{enumerate}
 \item $11 + 5 \times 4$
 \item $ 8 \div 2 \times 2$
 \item $ 12 \div 2 + 4$
 \item $ (23 - 1 + 4) \div 2 + 7$
 \item $ 4 + 2 \times 3 + 4$
 \end{enumerate} 
 \subsection*{Exercice 2 (3 points)}
 
 \begin{enumerate}
 \item Donner la fraction dont le dénominateur est $7$ et le numérateur est $12$. 
 \item Quel est le numérateur de $\frac73$ ?
 \item Donner une fraction dont le dénominateur est le double du numérateur. 
 \end{enumerate}
 
 \subsection*{Exercice 3 (6 points)}
Trouver l'écriture décimale des fractions suivantes (si elle est infinie, entourez la période se répétant): 
\begin{enumerate}
\item $\frac45$
\item $\frac38$
\item $\frac43$
\item $\frac{13}{40}$
\item $\frac37$
\item $\frac{497}{100}$
\end{enumerate}
 \subsection*{Exercice 4 (6 points)}
 
 Simplifier les fractions suivantes : 
 \begin{enumerate}
 \item $\frac{19}{38}$
 \item $\frac{15}{75}$
 \item $\frac{56}{49}$
 \item $\frac{495}{330}$
 \end{enumerate}
 
 
 
\newpage  

 \subsection*{Exercice 1 (5 points)}
 
 Effectuer les calculs suivants en détaillant les étapes : 
 \begin{enumerate}
 \item $7 + 5 \times 4$
 \item $ 12 \div 2 \times 2$
 \item $ 24 \div 2 + 4$
 \item $ (5 - 1 + 4) \div 2 + 7$
 \item $ 4 + 2 \times 3 + 4$
 \end{enumerate} 
 \subsection*{Exercice 2 (3 points)}
 
 \begin{enumerate}
 \item Donner la fraction dont le numérateur est $14$ et le dénominateur est $9$. 
 \item Quel est le dénominateur de $\frac76$ ?
 \item Donner une fraction dont le dénominateur est le triple du numérateur. 
 \end{enumerate}
 
 \subsection*{Exercice 3 (6 points)}
Trouver l'écriture décimale des fractions suivantes (si elle est infinie, entourez la période se répétant): 
\begin{enumerate}
\item $\frac35$
\item $\frac58$
\item $\frac{10}3$
\item $\frac{19}{40}$
\item $\frac47$
\item $\frac{932}{100}$
\end{enumerate}
 \subsection*{Exercice 4 (6 points)}
 
 Simplifier les fractions suivantes : 
 \begin{enumerate}
 \item $\frac{13}{26}$
 \item $\frac{135}{75}$
 \item $\frac{35}{49}$
 \item $\frac{495}{825}$
 \end{enumerate}
 
 
 
\newpage  
 
 	\end{document}
 	
